\documentclass[12pt,oneside]{book}

% ===================== METADATOS =====================
\newcommand{\universidad}{Universidad de Buenos Aires}
\newcommand{\facultad}{Facultad de Derecho}
\newcommand{\materia}{Derecho Procesal Civil y Comercial}
\newcommand{\titulo}{Notificaciones, Rebeldía y Práctica Profesional en el Proceso de Familia}
\newcommand{\autor}{Isaac Edgar Camacho Ocampo}
\newcommand{\anio}{2026}

% ===================== PAQUETES =====================
\usepackage[utf8]{inputenc}
\usepackage[T1]{fontenc}
\usepackage[spanish,es-nodecimaldot]{babel}
\usepackage[top=2.5cm,bottom=2.5cm,left=3cm,right=2.5cm,headheight=15pt]{geometry}
\usepackage{amsmath}
\usepackage{amssymb}
\usepackage{mathtools}
\usepackage{graphicx}
\usepackage{float}
\usepackage{caption}
\usepackage{subcaption}
\usepackage{booktabs}
\usepackage{array}
\usepackage{longtable}
\usepackage{tabularx}
\usepackage{enumitem}
\usepackage{xcolor}
\usepackage[most]{tcolorbox}
\usepackage{fancyhdr}
\usepackage{ragged2e}
\usepackage[
    colorlinks=true,
    linkcolor=blue!50!black,
    citecolor=green!40!black,
    urlcolor=blue!60!black,
    pdftitle={\titulo},
    pdfauthor={\autor},
    pdfsubject={Apuntes universitarios},
    bookmarks=true
]{hyperref}

% ===================== CONFIGURACIÓN =====================
\graphicspath{{./img/}{./graficos/}{./figuras/}}
\setcounter{tocdepth}{2}
\setlength{\parindent}{0pt}
\setlength{\parskip}{0.5em}
\setlist[itemize]{topsep=0.3em,itemsep=0.2em}
\setlist[enumerate]{topsep=0.3em,itemsep=0.2em}
\clubpenalty=10000
\widowpenalty=10000

\pagestyle{fancy}
\fancyhf{}
\fancyhead[L]{\nouppercase{\leftmark}}
\fancyhead[R]{\materia}
\fancyfoot[C]{\thepage}
\renewcommand{\headrulewidth}{0.4pt}
\renewcommand{\footrulewidth}{0pt}
\fancypagestyle{plain}{
    \fancyhf{}
    \fancyfoot[C]{\thepage}
    \renewcommand{\headrulewidth}{0pt}
}

% ===================== COMANDOS =====================
\newcommand{\abs}[1]{\left|#1\right|}
\newcommand{\norm}[1]{\left\lVert#1\right\rVert}
\newcommand{\vect}[1]{\mathbf{#1}}
\newcommand{\deriv}[2]{\frac{d#1}{d#2}}
\newcommand{\pderiv}[2]{\frac{\partial#1}{\partial#2}}

% ===================== CAJAS =====================
\newtcolorbox{definition}[1]{enhanced,breakable,title={Definición: #1},fonttitle=\bfseries,colback=blue!3,colframe=blue!50!black,boxrule=0.8pt,arc=2mm}
\newtcolorbox{important}{enhanced,breakable,title={Importante},fonttitle=\bfseries,colback=yellow!8,colframe=orange!70!black,boxrule=0.8pt,arc=2mm}
\newtcolorbox{example}[1]{enhanced,breakable,title={Ejemplo: #1},fonttitle=\bfseries,colback=green!4,colframe=green!40!black,boxrule=0.8pt,arc=2mm}
\newtcolorbox{formula}[1]{enhanced,breakable,title={Fórmula / Regla: #1},fonttitle=\bfseries,colback=purple!4,colframe=purple!50!black,boxrule=0.8pt,arc=2mm}
\newtcolorbox{exercise}[1]{enhanced,breakable,title={Ejercicio: #1},fonttitle=\bfseries,colback=gray!5,colframe=gray!60!black,boxrule=0.8pt,arc=2mm}
\newtcolorbox{note}{enhanced,breakable,title={Nota},fonttitle=\bfseries,colback=white,colframe=gray!50,boxrule=0.6pt,arc=2mm}

\begin{document}

% ===================== PORTADA =====================
\begin{titlepage}
\thispagestyle{empty}
\begin{center}
\vspace*{1cm}
{\large \textsc{\universidad}}

\vspace{0.2cm}
{\large \textsc{\facultad}}

\vspace{3cm}
{\Huge \textbf{\materia}}

\vspace{1cm}
{\LARGE \titulo}

\vspace{0.8cm}
\textit{Comprender los conceptos es el primer paso para convertir el estudio en conocimiento.}

\vspace{1.2cm}
\rule{0.70\textwidth}{0.5pt}

\vspace{0.5cm}
{\large Apuntes de estudio}

\vspace{0.5cm}
\rule{0.70\textwidth}{0.5pt}

\vfill
{\large \autor}

\vspace{0.3cm}
{\large \anio}
\end{center}
\vfill
\begin{flushleft}
\textit{Este es un modesto aporte a los estudiantes de ``Derecho Procesal Civil y Comercial'' no pretende sustituir el material oficial de la materia.}
\end{flushleft}
\end{titlepage}

% ===================== ÍNDICE =====================
\tableofcontents

% ===================== CONTENIDO =====================
\chapter[Demanda, preclusión y notificaciones]{Marco teórico procesal: la demanda, el traslado y las notificaciones}

\section{El trámite inicial de la demanda}

El circuito de la demanda se inicia con su presentación. Una vez que la demanda está presentada y completa en el juzgado, el juez dicta el \textbf{primer proveído}. Como regla de control de plazos, si el estado del expediente no se verifica dentro de las veinticuatro horas, corresponde dar aviso; el docente señaló este punto como casi lo más importante del circuito.

\subsection{El primer proveído y el traslado}

En el primer proveído el juez determina qué hay, qué no hay, qué falta y qué no falta. Si no hay nada que deba aclararse o exigirse, ordena el \textbf{traslado}. Cuando el juez verifica que las partes, el objeto, el hecho y el derecho están correctamente expuestos, ordena el traslado.

\begin{definition}{Proveído con ``previo''}
Cuando un proveído dice \textbf{``previo''}, significa que antes de dar traslado a la otra parte es necesario cumplir algo; si ello no se cumple, el trámite no avanza.
\end{definition}

\section{El principio de preclusión}

El proceso civil, como los procesos del sistema jurídico en general, es \textbf{``estanco''}: una vez cerrada una etapa y pasada a la siguiente, no se vuelve atrás. Lo que no se hizo en su momento no puede realizarse después, salvo excepciones. Si la contraparte está atenta, hará saber que la etapa ya pasó y que se perdió un derecho.

\begin{definition}{Principio de preclusión}
Principio según el cual el proceso avanza hacia adelante y no retrocede: una vez superada una etapa, no se vuelve a ella, y lo que no se hizo en la etapa correspondiente se pierde como derecho.
\end{definition}

% TODO: contenido incompleto o ambiguo en las notas originales: el docente mencionó una excepción ("excepto...") sin precisarla.

\section{Las notificaciones y el artículo 135 del CPCCN}

\subsection{Lectura del artículo 135}

Se asignó como lectura el \textbf{artículo 135} del Código Procesal Civil y Comercial de la Nación (CPCCN), que establece las formas de notificación, con atención especial a las \textbf{cédulas}, que son varias. Según el docente, la norma no dice mucho más porque el Código es muy antiguo. En la clase siguiente se verá la diferencia entre notificación electrónica y notificación en papel.

% TODO: contenido incompleto o ambiguo en las notas originales: el apunte no detalla el contenido del art. 135; el documento Word incluía un resumen de la norma elaborado fuera de la clase, que no se reproduce.

\subsection{Uso de códigos comentados}

\begin{important}
El docente exigió contar con un \textbf{código comentado}, tanto del Código Civil y Comercial como del Código Procesal, en el formato que cada estudiante prefiera. Indicó que la biblioteca digital de la facultad da acceso a dos bases de datos jurídicas y que en alguna de ellas figura el comentario del artículo.
\end{important}

\section{La cédula electrónica en el expediente digital}

Para la próxima clase se indicó leer el artículo 135 y observar qué corresponde a notificación electrónica y qué a notificación en papel. La cédula electrónica puede verificarse en el \textbf{PJN}: en cada expediente, en todos los recién iniciados, hay una cédula electrónica.

\section{Cuadro Cornell: demanda, preclusión y notificaciones}

\begin{table}[H]
\centering
\begin{tabularx}{\textwidth}{
    >{\raggedright\arraybackslash}p{0.27\textwidth}
    >{\raggedright\arraybackslash}X
}
\toprule
\textbf{Preguntas / palabras clave} & \textbf{Notas / respuestas} \\
\midrule
\textbf{¿Qué hace el juez al recibir la demanda y qué ordena el primer proveído?} &
Dicta el primer proveído, en el que determina qué hay, qué no hay, qué falta y qué no falta. Si todo está correcto (partes, objeto, hecho y derecho), ordena el traslado. \\[0.6em]

\textbf{¿Qué significa que un proveído diga ``previo''?} &
Que antes de pasar a la otra parte debe cumplirse algo; si no se cumple, el trámite no avanza. \\[0.6em]

\textbf{¿Qué es el principio de preclusión?} &
El proceso es estanco: una vez superada una etapa no se vuelve a ella, y lo que no se hizo en su momento se pierde como derecho. Va hacia adelante y no retrocede. \\[0.6em]

\textbf{¿Qué se debe leer del artículo 135 del CPCCN y por qué?} &
Establece las formas de notificación, en especial las cédulas, que son varias. Debe leerse con un código comentado. En la clase siguiente se distinguirá la notificación electrónica de la de papel. \\[0.6em]

\textbf{¿Dónde se verifica una cédula electrónica?} &
En el PJN: en cada expediente, en todos los recién iniciados, hay una cédula electrónica. \\
\midrule
\multicolumn{2}{p{\dimexpr\textwidth-2\tabcolsep\relax}}{\textbf{Resumen:} tras la presentación de la demanda, el juez dicta el primer proveído y, si todo está en orden, ordena el traslado; un ``previo'' exige cumplir algo antes de continuar. El proceso se rige por la preclusión: avanza sin retroceder. Las notificaciones se estudian a partir del artículo 135 del CPCCN, con códigos comentados, y la cédula electrónica puede verificarse en el PJN.} \\
\bottomrule
\end{tabularx}
\end{table}

% ---------------------------------------------------------
\chapter[Rebeldía y proceso de alimentos]{La rebeldía y su excepción en el derecho de familia}

\section{La declaración de rebeldía (artículo 59 del CPCCN)}

Se asignó como lectura el \textbf{artículo 59} del CPCCN, descripto por el docente como simple y concreto. La regla básica expuesta es: quien no contesta cuando debe contestar, es declarado rebelde. La rebeldía se declara mediante un proveído o resolución del juez.

\begin{definition}{Rebeldía}
Situación en la que, según lo expuesto en clase, es declarada la parte que no contesta cuando tiene que contestar. Se declara por resolución del juez (artículo 59 del CPCCN).
\end{definition}

% TODO: contenido incompleto o ambiguo en las notas originales: el docente preguntó qué debe hacerse con la resolución que declara la rebeldía, pero la respuesta no quedó registrada. El texto del art. 59 incluido en el Word fue una reconstrucción ajena a la clase y no se reproduce.

\section{La rebeldía en el proceso de alimentos}

\begin{important}
\textbf{En alimentos no hay rebeldía.} Pedir la rebeldía cuando el alimentante no contesta la demanda es, según el docente, un error grave.
\end{important}

El fundamento es práctico: el proceso de familia de alimentos busca que el alimentante \textbf{pague siempre}. Declarar la rebeldía implica sacar a la parte del proceso por un tiempo, y el proceso de alimentos nunca quiere sacarlo, sino mantenerlo dentro para que pague y se haga cargo de los alimentados. Además, en cualquier momento del proceso el alimentante puede ingresar al proceso de alimentos de la manera que fuera.

\section{Cuadro Cornell: rebeldía}

\begin{table}[H]
\centering
\begin{tabularx}{\textwidth}{
    >{\raggedright\arraybackslash}p{0.27\textwidth}
    >{\raggedright\arraybackslash}X
}
\toprule
\textbf{Preguntas / palabras clave} & \textbf{Notas / respuestas} \\
\midrule
\textbf{¿Qué ocurre si la parte demandada no contesta cuando debe hacerlo?} &
Es declarada rebelde (artículo 59 del CPCCN). \\[0.6em]

\textbf{¿Cómo se declara la rebeldía?} &
Mediante un proveído o resolución del juez. \\[0.6em]

\textbf{¿Procede la rebeldía en los procesos de alimentos?} &
No. Pedir la rebeldía si el alimentante no contesta la demanda es un error grave. \\[0.6em]

\textbf{¿Por qué no hay rebeldía en alimentos?} &
Porque el proceso busca que el alimentante pague siempre: la rebeldía lo sacaría del proceso por un tiempo, y el proceso quiere mantenerlo dentro para que se haga cargo de los alimentados. Además, el alimentante puede ingresar al proceso en cualquier momento. \\
\midrule
\multicolumn{2}{p{\dimexpr\textwidth-2\tabcolsep\relax}}{\textbf{Resumen:} la regla del artículo 59 del CPCCN es que quien no contesta cuando debe es declarado rebelde por resolución judicial. En alimentos no corresponde pedirla, porque el proceso procura mantener al alimentante dentro de la causa para que pague y se haga cargo de los alimentados.} \\
\bottomrule
\end{tabularx}
\end{table}

% ---------------------------------------------------------
\chapter[Taller de redacción]{Taller práctico de redacción: medida de vinculación provisoria}

Se trabajó en forma colectiva y en tiempo real, sobre una plataforma de edición compartida, en la corrección de un escrito de un proceso de familia que solicita una \textbf{vinculación provisoria}. El pedido se vincula con el temor manifestado por una de las partes.

\section{Domicilio real en el escrito}

Al escrito debía incorporarse el domicilio real omitido. Se advirtió que debía modificarse también el encabezado. Al cargar la dirección se corrigió la ciudad: no bastaba el dato de la calle y número, sino que había que aclarar la localidad (San Miguel, provincia de Buenos Aires), pues una búsqueda arrojaba una localidad homónima distinta. La sección se tituló \textbf{``Se constituye domicilio real''}.

% TODO: contenido incompleto o ambiguo en las notas originales: el apunte usa "constituye domicilio real"; el Word incluía una distinción entre domicilio real y constituido que no fue expuesta en clase y no se reproduce.

\section{El petitorio}

No bastaba agregar el domicilio: debía \textbf{modificarse el petitorio}. En la redacción colectiva se descartó la fórmula ``tenga'' y se acordó una fórmula del tipo \textbf{``se ordene la vinculación provisoria solicitada en el punto\ldots''}, ubicada en el \textbf{punto cuatro} del petitorio. Se aclaró que lo relativo al domicilio corresponde a ``se constituye domicilio real'' y no al pedido de vinculación.

\begin{note}
La parte final del taller consistió en la coordinación de la firma y la carga del escrito en el sistema, con indicación de que debía subirse el mismo día.
\end{note}

\section{Cuadro Cornell: taller de redacción}

\begin{table}[H]
\centering
\begin{tabularx}{\textwidth}{
    >{\raggedright\arraybackslash}p{0.27\textwidth}
    >{\raggedright\arraybackslash}X
}
\toprule
\textbf{Preguntas / palabras clave} & \textbf{Notas / respuestas} \\
\midrule
\textbf{¿Qué se solicitaba en el escrito del taller y por qué?} &
Una vinculación provisoria, vinculada al temor manifestado por una de las partes. \\[0.6em]

\textbf{¿Qué se corrigió en cuanto al domicilio?} &
Se incorporó el domicilio real omitido, se aclaró la ciudad (San Miguel, provincia de Buenos Aires) y se tituló la sección ``Se constituye domicilio real''. También debía ajustarse el encabezado. \\[0.6em]

\textbf{¿Cómo se incorpora un pedido nuevo al petitorio?} &
Modificando el petitorio: se redacta el pedido y se lo ubica en su punto (aquí, el punto cuatro) con la fórmula ``se ordene la vinculación provisoria solicitada''. \\
\midrule
\multicolumn{2}{p{\dimexpr\textwidth-2\tabcolsep\relax}}{\textbf{Resumen:} en el taller se corrigió un escrito de familia que pide una vinculación provisoria: se incorporó y tituló el domicilio real con su ciudad, se ajustó el encabezado y se modificó el petitorio con un punto específico para ese pedido.} \\
\bottomrule
\end{tabularx}
\end{table}

\end{document}
